% AI-produced by GPT-6 Astra (OpenAI, ChatGPT Pro mode). No priority claim.
\subsection{The binary cycle obstruction vanishes in the full category}

We retain either category $A$ of
Proposition~\ref{A:predimension-quasi}. Let $G_X$ have vertices
$A(V,X)$ and an edge $e:W\to X$ directed from $ev_1$ to $ev_2$.
Put
\[
 K(X)=\ker\bigl(\Z[A(W,X)]\xrightarrow{v_2-v_1}
                         \Z[A(V,X)]\bigr).
\]
Thus $K(X)$ is the integral cycle group of $G_X$. The augmentation
$\epsilon:K(X)\to\Z$ sums its edge coefficients.

\begin{lemma}[Proved: adjoining a loose path of length three]
\label{A:predimension-three-path}
Let $B$ be a positive-predimension hypergraph and let $x,c$ be two
distinct vertices of $B$. Adjoin five distinct new vertices
$r,s,a,b,d$ and precisely the three new hyperedges
\[
 \{x,r,a\},\qquad \{r,s,b\},\qquad \{s,c,d\}.
 \tag{LP}\label{eq:predimension-loose-path}
\]
Then $B$ is $d$-closed in the resulting positive-predimension
hypergraph. No assumption that $\{x,c\}$ is $d$-closed is needed.
\end{lemma}
\begin{proof}
For a nonempty subset $N$ of the five new vertices, let $e(N)$ be
the number of the three displayed edges whose new vertices all lie
in $N$. The relative predimension over the whole of $B$ is
$2|N|-3e(N)$. If $e(N)=0$, this is positive. If $e(N)=1$, then
$|N|\geq2$, so it is at least one. If $e(N)=2$, the union of the
new vertices required by those edges has size at least four, so it
is at least two. If $e(N)=3$, all five new vertices occur and it is
one. Thus every proper extension of $B$ has strictly larger
predimension.

For an arbitrary subset $U$ meeting the new vertices,
submodularity gives
\[
 \delta(U)-\delta(U\cap B)
 \geq\delta(U\cup B)-\delta(B)>0.
\]
Since the old subsets have nonnegative predimension, this proves
positivity of every nonempty subset of the extension.
\end{proof}

\begin{lemma}[Proved: atoms in a loose cycle]
\label{A:predimension-loose-cycle}
Let $C$ be a loose hypergraph cycle of length $\ell\geq4$: it has
binary vertices $b_0,\ldots,b_{\ell-1}$, distinct private vertices
$p_0,\ldots,p_{\ell-1}$, and edges
$\{b_i,b_{i+1},p_i\}$ with cyclic indices. Then $C$ has positive
predimension, and every point and every hyperedge is $d$-closed.
\end{lemma}
\begin{proof}
For a chosen subset $B'$ of the binary vertices, including a private
vertex lowers predimension only when both its binary endpoints are
included; in that case its net contribution is $-1$. Otherwise its
contribution is $2$. Thus the minimum at a fixed nonempty $B'$ is
\[
 2|B'|-e(B'),
\]
where $e(B')$ counts the edges of the ordinary cycle induced on
$B'$. A proper nonempty $B'$ induces a forest, so this minimum is
at least $|B'|+1$; the whole set gives $\ell$. With no binary
vertices, every nonempty subset has predimension at least two.
This proves positivity.

A proper extension of a hyperedge either adds a private vertex
with its binary endpoints unchanged, increasing predimension by
two, or has at least three binary vertices. In the latter case the
preceding bound is at least four, including the full-cycle case
$\ell\geq4$. Its predimension therefore exceeds the edge's value
three. Every nonempty subset of at least two vertices has
predimension at least three: with two or more binary vertices use
the same bound, and with at most one binary vertex it has no edge.
Thus every proper extension of a point exceeds its predimension
two as well.
\end{proof}

\begin{proposition}[Proved: the binary cycle colimit is exactly its augmentation]
\label{A:predimension-binary-colimit}
For either category $A$ of
Proposition~\ref{A:predimension-quasi}, augmentation induces an
isomorphism
\[
 \operatorname{colim}_A K\xrightarrow{\ \epsilon\ }\Z.
\]
Consequently, for
$H=A(V,-)/(ev_1\sim ev_2)$,
\[
 H_2(B\operatorname{El}_A H;\Z)=0.
\]
In particular, the binary class used to prove nontotality for the
selected generated AKR category vanishes upon passage to this full
predimension category.
\end{proposition}
\begin{proof}
Fix $z=\sum_e c_e e\in K(X)$, with nonempty finite edge support;
the zero cycle needs no construction. Let $B$ be the set of its
binary endpoints. These are actual $V$-arrows and hence $d$-closed
points of $X$. Adjoin a new centre $c$, and for each $x\in B$ a
loose length-three path
\[
 P_x=(x\longrightarrow r_x\longrightarrow s_x\longrightarrow c)
\]
using five new vertices $r_x,s_x,a_x,b_x,d_x$ as in
\eqref{eq:predimension-loose-path}. All these new vertices are
distinct for distinct $x$. Call the resulting hypergraph $Z$.
Starting from the disjoint union of $X$ and $\{c\}$, repeated
applications of Lemma~\ref{A:predimension-three-path} prove that
$Z$ is positive-predimension and $X\leq_d Z$.

For a support edge $e$ write its ordered vertices as $u<v<p$.
Its binary endpoints are $u,v$. Let
\[
 F_e=e\cup P_u\cup P_v\subseteq Z.
\]
This is a loose seven-cycle, containing its old edge as a
$d$-closed subobject by Lemma~\ref{A:predimension-loose-cycle}.
The old edge is also $d$-closed in $X$, since $e$ is an actual
$W$-arrow. The induced subhypergraph $X\cup F_e$ is precisely the
free amalgam of $X$ and $F_e$ over that edge. Thus
$F_e\leq_d X\cup F_e$. Adjoin all remaining paths $P_x$ one at
a time. Their endpoints $x,c$ already exist, and
Lemma~\ref{A:predimension-three-path} makes every step $d$-closed.
Transitivity gives $F_e\leq_d Z$.

This proof allows $p$ to be a binary endpoint of another support
edge. In that case its remaining path joins two existing points
of $F_e$ and still has strictly positive relative predimension.
This is why paths of length three are used.

Extend the order of $X$ as follows. Place all old points first.
Next place the two-point blocks $r_x<s_x$, in increasing order of
$x\in B$. Place $c$ after all those blocks. Finally place the
three-point blocks $a_x<b_x<d_x$, again in increasing order of
$x$. All new hyperedges then have their private point last, so
$P_x$ is an increasing directed binary path. The order induced on
every $F_e$ has the same pattern:
\[
 u<v<p<r_u<s_u<r_v<s_v<c
       <a_u<b_u<d_u<a_v<b_v<d_v.
\]
Hence all $F_e$ are copies, by increasing induced embeddings, of
one fixed ordered seven-cycle object $P$. Its distinguished cycle
\[
 \kappa=e+P_v-P_u
\]
has augmentation $1+3-3=1$. Every point and edge of $P$ is
$d$-closed by Lemma~\ref{A:predimension-loose-cycle}. The maps
$P\cong F_e\hookrightarrow Z$ are therefore arrows in $A$.
If $A$ is the point-and-edge-closed subcategory, $Z$ also lies in
it: its old atoms are closed by $X\leq_d Z$, and every new atom
lies in some $F_e\leq_d Z$, because each $x\in B$ is incident
with a support edge.

Let $z_e=e+P_v-P_u\in K(Z)$ be the image of $\kappa$ for edge $e$.
The coefficient of $P_x$ in $\sum_e c_e z_e-z$ is the coefficient
of $x$ in $\partial z$, which is zero. Thus in the actual cycle
group of the one target $Z$,
\[
 z=\sum_e c_e z_e.
\]
Passing to the colimit yields
\[
 [z]=\left(\sum_e c_e\right)[\kappa]=\epsilon(z)[\kappa].
\]
This holds for every cycle at every object. Since $\epsilon(\kappa)=1$,
it proves that augmentation on the colimit is an isomorphism.

Finally, the exact sequence
\[
 0\to K\to\Z[A(W,-)]\xrightarrow{v_2-v_1}
 \Z[A(V,-)]\to\Z[H]\to0
\]
and projectivity of the two representable modules give
\[
 H_2(B\operatorname{El}_A H;\Z)
 \cong\ker\bigl(\operatorname{colim}_A K
                         \xrightarrow{\epsilon}\Z\bigr)=0.
\]
This is the same natural binary comparison as in
Proposition~\ref{C:attachment-cycle-test}.
\end{proof}

\paragraph{Scope of the binary calculation.}
The cone calculation proves vanishing for this one binary action in
degree two. Totality requires the separate all-action argument below;
the all-ring triple obstruction of
Proposition~\ref{A:predimension-triple} remains nonzero.
The construction also works for cycles with coefficients in any
commutative ring, giving the corresponding degree-two vanishing.
